\begingroup\fontsize{10}{11.8}\selectfont\setlength{\abovedisplayskip}{5pt}\setlength{\belowdisplayskip}{5pt}
\subsection{Calendar census and simultaneous compatibility of the regional biographies}
The preceding rows are assembled, without yet applying any decimal reader,
into the three block biographies produced by the transitions:
\begin{align}
 B^{\rm cl}
  &=010211|012222|010211|002111|110221,\notag\\
 B^{\rm pr}
  &=201101|121221|102011|012222|102011,
 \label{act21:eq:biografias-tpk-tres-caracteres}\\
 B^{\rm au}
  &=121200|112202|121020|010210|010200.\notag
\end{align}
Indeed, the consecutive pairs of the first two segments are the rows
of \(X_0,Y_0\), and those of the last two are the rows of \(X_1,Y_1\).
Thus, the biographies are another representation of the record of twelve
transitions, not a collection of conventional prefixes.

To verify the operative word of the calendar, let
\(c=c_1c_2c_3c_4\in\{0,1\}^4\) and, for each regime
\(\chi\in\{{\rm cl},{\rm pr},{\rm au}\}\), define
\[
 u_0^\chi=b_0^\chi,\qquad
 u_j^\chi=u_{j-1}^\chi L_{c_j}\quad(1\le j\le4).
\]
To condense the census, denote by
\[
 N_{\mathrm{ar}}(c):=
 \#\{(j,\chi):u_j^\chi=b_j^\chi,\ 1\le j\le4\},\qquad
 N_{\mathrm{bio}}(c):=
 \#\{\chi:(u_0^\chi,\ldots,u_4^\chi)=B^\chi\}.
\]
Exact multiplication in \(\FF_3\) for the sixteen calendars gives:
\[
\begin{array}{c|c|c}
 c&N_{\mathrm{ar}}(c)&N_{\mathrm{bio}}(c)\\ \hline
0000,0001&6&0\\
0010&9&0\\
0011&12&3\\
0100,0101,0110,0111&3&0\\
1000,1001,1010,1011&0&0\\
1100,1101,1110,1111&0&0
\end{array}
\tag{D}
\label{act21:eq:censo-dieciseis-calendarios}
\]
The table exhausts the \(2^4=16\) possibilities, and each of its entries is
obtained by four multiplications of a row by one of the two displayed matrices.
Consequently, \(0011\) is the unique calendar that simultaneously reconstructs
the three complete biographies. In operator notation,
\begin{equation}
 (L_0,L_0,L_1,L_1)
 \label{act21:eq:calendario-0011-interno}
\end{equation}
is not an independent prescription: it is the unique binary reading compatible
with the twelve edges already produced by \(U_t\).

The calendar \eqref{act21:eq:calendario-0011-interno} produces four new
six-component blocks. Their concatenations form
\[
 w_6\longrightarrow w_{12}\longrightarrow w_{18}
 \longrightarrow w_{24}\longrightarrow w_{30}.
\]
The subscripts of \(w\) indicate cumulative length. The boundary \(R_{36}\) preserves the terminal structure and opens the prolongation relation; it is not renamed as a final periodic block that could be repeated indefinitely.

The identity \(L=X^{-1}Y\) proves uniqueness once the transitions are
fixed; it does not replace their production by \(U_t\). Conversely,
\eqref{act21:eq:censo-dieciseis-calendarios} proves within the article the uniqueness
of the ordering of the four transports. The causal chain is thereby closed:
the APP census produces the initial regions, TRIT fixes regime and
orientation, TPK produces the transition record, and that record
determines \(L_0,L_1\) and \(0011\) before any Archimedean realization.


\par\endgroup
